\section{Consequences for heat-bath dynamics}\label{sec:dynamics}

The preceding section provides the static inputs used here. A top-eigenvector
projection rarely lies near zero, while $Q_n=O_{\Prob}(n^{-2/3})$ and the
covariance norm has the critical scale. A stationary increment estimate
converts the projection bound into a sign test for each realized start.
For linear functions, a cavity comparison uses the squared-overlap bound
to control the contribution of sites with small average conditional variance,
and then compares inverse Rayleigh quotients with the covariance norm.
Finally, a density comparison with one fixed subcritical temperature proves
the subexponential upper bound and completes Theorem~\ref{thm:mixing}.

\subsection{Stationary increments and typical starting configurations}

\begin{lemma}[Stationary increments]\label{lem:stationary-increments}
Fix the disorder and a function $f:\{-1,1\}^n\to\R$.
Let $(X_t)_{t\ge0}$ be the continuous-time heat-bath chain with generator
$\mathcal L$, and let $(Y_k)_{k\ge0}$ be the discrete-time chain with
transition operator $P$, each started from $\mu$.  Then
\begin{align}
 \E_\mu\bigl[(f(X_t)-f(X_0))^2\bigr]
 &\le 2t\,\mathcal D_\mu(f), \qquad t\ge0,\label{eq:stationary-increments-ct}\\
 \E_\mu\bigl[(f(Y_k)-f(Y_0))^2\bigr]
 &\le \frac{2k}{n}\,\mathcal D_\mu(f),
 \qquad k\in\{0,1,2,\ldots\}.\label{eq:stationary-increments-dt}
\end{align}
Here the expectations average only over the initial configuration and
the trajectory.  In particular, if $g(x)=v^Tx$ and $\|v\|=1$, the two
right-hand sides are at most $2t$ and $2k/n$, respectively.
\end{lemma}

\begin{proof}
For either clock, reversibility diagonalizes the mean-square
increment in the same basis as the Dirichlet form. In the inner product
$\langle f,h\rangle_\mu=\mu[fh]$, stationarity gives
\[
 \E_\mu\bigl[(f(X_t)-f(X_0))^2\bigr]
 =2\langle f,(\mathrm{Id}-e^{t\mathcal L})f\rangle_\mu.
\]
The operator $-\mathcal L$ is self-adjoint and nonnegative.  Expanding in
an orthonormal eigenbasis and using $1-e^{-t\lambda}\le t\lambda$ gives
\eqref{eq:stationary-increments-ct}, because
$\langle f,-\mathcal Lf\rangle_\mu=\mathcal D_\mu(f)$ by
Lemma~\ref{lem:heatbath-form}.

Each conditional expectation $P_i$ is an orthogonal projection in
$L^2(\mu)$.  Thus $P=n^{-1}\sum_iP_i$ is self-adjoint with spectrum in
$[0,1]$.  The discrete stationary increment is
$2\langle f,(\mathrm{Id}-P^k)f\rangle_\mu$.
For $0\le\theta\le1$, $1-\theta^k\le k(1-\theta)$, while
$\langle f,(\mathrm{Id}-P)f\rangle_\mu=\mathcal D_\mu(f)/n$.
This proves \eqref{eq:stationary-increments-dt}.
Finally, Lemma~\ref{lem:heatbath-form} gives
\[
 \mathcal D_\mu(v^Tx)
 =\sum_{i=1}^n v_i^2\mu[\Var(x_i\mid x_{-i})]
 \le\sum_{i=1}^n v_i^2=1.
\]
\end{proof}

\begin{proof}[Proof of Theorem~\ref{thm:gibbs-start}]
We prove the continuous-time statement in three steps. Choose a
projection threshold above its root-mean-square increment scale, use its
sign as a test for each fixed starting state, and average only the
resulting intrinsic bad-state indicator. This last step will remove the
auxiliary diagonal without averaging away the distance being tested.

Let $t_n=o(n^{2/3})$ be deterministic and nonnegative. Choose a positive
deterministic $b_n\to0$ such that
\[
 \frac{t_n}{b_n^2n^{2/3}}\longrightarrow0;
\]
for example, $b_n=(t_n/n^{2/3})^{1/4}+n^{-1}$ works.  Put
$a_n=b_nn^{1/3}$.

\smallskip\noindent\emph{A sign that persists.}
Separate the auxiliary randomness by writing $W=A+D$, where $A$
has zero diagonal and $D$ is the independent Gaussian diagonal.  The
Gibbs law $\mu_A=\mu_W$ and the heat-bath transition laws depend only on
$A$.  The unit top eigenvector $v=v(A+D)$ is allowed to depend on both
matrices.  Set $g(x)=v^Tx$.  Conditional on $A,D$, start the trajectory
from $\mu_A$.  If $|g(X_0)|>a_n$, failure of the two endpoint projections
to have the same nonzero sign requires
$|g(X_{t_n})-g(X_0)|\ge a_n$.  Hence
Lemma~\ref{lem:stationary-increments} and Markov's inequality give
\begin{equation}\label{eq:sign-failure}
 \Prob_{\mu_A}\bigl(g(X_0),g(X_{t_n})
     \text{ do not have the same nonzero sign}\mid A,D\bigr)
 \le \mu_A\{|g|\le a_n\}+\frac{2t_n}{a_n^2}.
\end{equation}
By Theorem~\ref{thm:static-small-ball}, the first term tends to zero in
probability over $A,D$.  It also tends to zero in expectation, since it
lies in $[0,1]$.  Therefore the deterministic quantity
\begin{equation}\label{eq:sign-error-average}
 \varepsilon_n
 :=\E_{A,D}\mu_A\{|g|\le a_n\}+\frac{2t_n}{a_n^2}
 \longrightarrow0
\end{equation}
bounds the average sign-failure probability.

\smallskip\noindent\emph{Testing each realized state.}
We now fix the initial configuration before defining its test event.
For $x$ with $g(x)\ne0$, let
\[
 H_{v,x}=\{y:g(y)g(x)>0\},\qquad
 r_n(A,D,x)=1-(e^{t_n\mathcal L}\mathbf1_{H_{v,x}})(x).
\]
For $g(x)=0$, set $r_n(A,D,x)=1$.  This definition includes every zero
projection in the failure event, and \eqref{eq:sign-failure} implies
\[
 \E_{A,D}\mu_A[r_n(A,D,\cdot)]\le\varepsilon_n.
\]
Spin-flip symmetry gives $\mu_A(H_{v,x})\le1/2$.  If $g(x)\ne0$,
testing total variation against this half-space yields
\[
 d_{t_n}^{\mathrm{ct}}(A,x)\ge\frac12-r_n(A,D,x).
\]
When $g(x)=0$ the following consequence remains valid because $r_n=1$:
\begin{equation}\label{eq:intrinsic-bad-state}
 \mathbf1_{\{d_{t_n}^{\mathrm{ct}}(A,x)\le1/4\}}
 \le4r_n(A,D,x).
\end{equation}
\smallskip\noindent\emph{Removing the auxiliary diagonal.}
The test half-space depends on $D$, but the bad-state indicator on the
left does not. We can therefore integrate \eqref{eq:intrinsic-bad-state}
to obtain
\[
 \E_A\mu_A\{x:d_{t_n}^{\mathrm{ct}}(A,x)\le1/4\}
 =\E_{A,D}\mu_A\{x:d_{t_n}^{\mathrm{ct}}(A,x)\le1/4\}
 \le4\varepsilon_n.
\]
In particular, for every fixed $\delta>0$,
\begin{equation}\label{eq:nested-gibbs-start}
 \Prob_A\!\left(
   \mu_A\{x:d_{t_n}^{\mathrm{ct}}(A,x)>1/4\}<1-\delta
 \right)
 \le\frac{4\varepsilon_n}{\delta}\longrightarrow0.
\end{equation}
This is the asserted convergence in probability over the original
off-diagonal disorder.  The argument integrates an intrinsic distance;
it does not require choosing one auxiliary diagonal that works for a
prescribed collection of initial states.

For deterministic integers $k_n=o(n^{5/3})$, use
$b_n=(k_n/n^{5/3})^{1/4}+n^{-1}$ and the discrete increment bound.
The same proof replaces $2t_n/a_n^2$ by $2k_n/(na_n^2)$ and proves
\eqref{eq:nested-gibbs-start} with $d_{k_n}^{\mathrm{dt}}$.
These are distances conditional on the realized initial configuration.
Averaging the transition law itself over a Gibbs initial configuration
would instead leave it equal to $\mu_A$ by stationarity.
\end{proof}

\subsection{Inverse Rayleigh quotients of linear functions}

The covariance estimate controls variances of linear functions.  To
control their inverse Rayleigh quotients, we must also account for sites
with small average conditional variance.

We use the standard spin-deletion cavity reweighting; see Talagrand
\cite{Talagrand2011MeanFieldI} and Chen \cite{Chen2013CavityLocal}.
Refined cavity and local-field limit theorems have been proved under
high-temperature hypotheses by Chatterjee \cite{Chatterjee2010Stein}
and Chen \cite{Chen2013CavityLocal}. The calculation below uses only
exact deletion identities and conditional Gaussian second moments, which
hold at the critical temperature.

\begin{proposition}[Linear tests]\label{prop:linear-scale}
With probability tending to one,
\begin{equation}\label{eq:linear-covariance-comparison}
 \|\Sigma\|\le\mathcal R_{\mathrm{lin}}(\mu)
 \le16\|\Sigma\|+n^{13/25}.
\end{equation}
In particular, $n^{-2/3}\mathcal R_{\mathrm{lin}}(\mu)$ is bounded in
probability.
\end{proposition}

\begin{proof}
The denominator in the inverse Rayleigh quotient of a linear function
is diagonal in its coefficients. We first rewrite the quotient as a
weighted covariance norm, then show that the sites with large weights
form a sufficiently small exceptional set.

\smallskip\noindent\emph{Weighted covariance.}
For each site define its local field and average conditional variance by
\[
 g_i(x)=\sum_{j\ne i}W_{ij}x_j,
 \qquad \alpha_i=\mu[\sech^2g_i].
\]
Thus $0<\alpha_i\le1$, and Lemma~\ref{lem:heatbath-form} gives
\[
 \mathcal D_\mu(a^Tx)=\sum_i\alpha_i a_i^2,
 \qquad \Var_\mu(a^Tx)=a^T\Sigma a.
\]
It follows by the change of variables $b_i=\sqrt{\alpha_i}a_i$ that
\begin{equation}\label{eq:weighted-covariance}
 \mathcal R_{\mathrm{lin}}(\mu)=\|\Gamma\|,
 \qquad
 \Gamma=\diag(\alpha_i^{-1/2})\Sigma\diag(\alpha_i^{-1/2}).
\end{equation}
For every nonzero $a$, the inequalities $\alpha_i\le1$ and the
nonnegativity of the variance give
\[
 \frac{a^T\Sigma a}{\sum_i\alpha_i a_i^2}
 \ge \frac{a^T\Sigma a}{\|a\|^2}.
\]
Taking the supremum proves the lower bound in
\eqref{eq:linear-covariance-comparison} deterministically.

\smallskip\noindent\emph{Deleting one spin.}
To control a site's weight, let $\nu_i$ be the Gibbs law on the other
$n-1$ spins after deleting
site $i$ and all its incident couplings, without rescaling the remaining
couplings.  Write $\bar\mu_i$ for the corresponding marginal of $\mu$ and
set
\[
 Z_i=\nu_i[\cosh g_i],\qquad
 R_i=\E_{\nu_i\otimes\nu_i}
       \left[\left(\frac{x_{-i}^Ty_{-i}}{n}\right)^2\right].
\]
Summing over $x_i$ gives the exact identities
\[
 \frac{d\bar\mu_i}{d\nu_i}=\frac{\cosh g_i}{Z_i},
 \qquad
 \alpha_i=\frac{\nu_i[(\cosh g_i)^{-1}]}{Z_i}
 \ge Z_i^{-2},
\]
where the last step is Jensen's inequality for the reciprocal function.
Since $\cosh g_i\ge1$, the product marginal comparison and
$(q-x_iy_i/n)^2\le2q^2+2n^{-2}$ imply
\begin{equation}\label{eq:cavity-overlap}
 R_i\le Z_i^2\E_{\bar\mu_i\otimes\bar\mu_i}
           \left[\left(\frac{x_{-i}^Ty_{-i}}n\right)^2\right]
 \le2Z_i^2(Q_n+n^{-2}).
\end{equation}

Let $\mathcal F_i$ be the sigma-field generated by the matrix with row
and column $i$ deleted.  Conditional on $\mathcal F_i$, the law $\nu_i$
is fixed and the incident couplings are independent $N(0,1/n)$
variables.  For a fixed spin configuration, $g_i$ is consequently a
centered Gaussian with variance $\sigma_{\mathrm{loc}}^2=(n-1)/n$.  For every fixed
integer $m\ge1$, Jensen's inequality yields
\[
 \E[Z_i^m\mid\mathcal F_i]
 \le \E_{G\sim N(0,\sigma_{\mathrm{loc}}^2)}[\cosh^mG]
 \le2e^{m^2/2}.
\]
Taking, for example, $m=200$ and using a union bound gives
\begin{equation}\label{eq:cavity-max-normalization}
 \Prob\left(\max_iZ_i>n^{1/100}\right)\longrightarrow0.
\end{equation}
Proposition~\ref{prop:overlap} says that $n^{2/3}Q_n$ is bounded in
probability.  Combining it with \eqref{eq:cavity-overlap} and
\eqref{eq:cavity-max-normalization}, we obtain
\begin{equation}\label{eq:cavity-small-overlaps}
 \Prob(E_n)\longrightarrow1,
 \qquad E_n=\left\{\max_iR_i\le r_n\right\},
 \qquad r_n=n^{-7/12}.
\end{equation}
For clarity, the exponent slack here is
$1/12-1/50=19/300>0$: on
$\max_iZ_i\le n^{1/100}$ and
$Q_n\le\tfrac14 n^{-2/3+19/300}$, the first term on the right of
\eqref{eq:cavity-overlap} is at most $r_n/2$, and its second term is
at most $r_n/2$ for all sufficiently large $n$.

\smallskip\noindent\emph{Counting exceptional sites.}
The uniform polynomial bound on $Z_i$ is not by itself enough. We also
need most $Z_i$ to stay below a fixed constant. Conditional on
$\mathcal F_i$, for two configurations of
the remaining spins put $c=x_{-i}^Ty_{-i}/n$.  Their two local fields
have variances $\sigma_{\mathrm{loc}}^2$ and covariance $c$, so
\[
 \E[\cosh g_i(x)\cosh g_i(y)\mid\mathcal F_i]
 =e^{\sigma_{\mathrm{loc}}^2}\cosh c.
\]
It follows that
\begin{align}
 \E[Z_i\mid\mathcal F_i]&=e^{\sigma_{\mathrm{loc}}^2/2},\label{eq:cavity-mean}\\
 \Var(Z_i\mid\mathcal F_i)
 &=e^{\sigma_{\mathrm{loc}}^2}\E_{\nu_i\otimes\nu_i}[\cosh c-1]
 \le C R_i,\label{eq:cavity-variance}
\end{align}
because $|c|\le1$ and $\cosh c-1\le Cc^2$ on this interval.

The event $E_n$ depends on many incident rows.  To use
\eqref{eq:cavity-variance} with its stated conditioning, introduce the
clipped count
\[
 N_n=\sum_{i=1}^n
       \mathbf1_{\{R_i\le r_n\}}\mathbf1_{\{Z_i>4\}}.
\]
The variable $R_i$ is $\mathcal F_i$-measurable, and
$4-e^{\sigma_{\mathrm{loc}}^2/2}\ge4-e^{1/2}>0$.  Thus conditional Chebyshev gives
\[
 \E N_n
 =\sum_i\E\!\left[
    \mathbf1_{\{R_i\le r_n\}}
    \Prob(Z_i>4\mid\mathcal F_i)\right]
 \le Cnr_n=Cn^{5/12}.
\]
On $E_n$, the clipped count equals $|B|$, where
$B=\{i:Z_i>4\}$.  Consequently,
\begin{equation}\label{eq:cavity-exceptional-count}
 \Prob(|B|>n^{1/2})
 \le\Prob(E_n^c)+\Prob(N_n>n^{1/2})
 \le\Prob(E_n^c)+Cn^{-1/12}\longrightarrow0.
\end{equation}
No independence between sites, or conditioning on $E_n$ in the Gaussian
calculation, is used.

\smallskip\noindent\emph{Combining the covariance blocks.}
The probabilistic work is complete. Work on the event
$\{\max_iZ_i\le n^{1/100},\ |B|\le n^{1/2}\}$, whose probability
tends to one by \eqref{eq:cavity-max-normalization} and
\eqref{eq:cavity-exceptional-count}.  Let $G=B^c$.
On $G$, $\alpha_i^{-1/2}\le Z_i\le4$, so the principal
blocks of the positive semidefinite matrix $\Gamma$ satisfy
\[
 \|\Gamma_{GG}\|\le16\|\Sigma\|,
 \qquad
 \|\Gamma_{BB}\|\le\Tr\Gamma_{BB}
 =\sum_{i\in B}\alpha_i^{-1}
 \le |B|\max_i Z_i^2\le n^{13/25}.
\]
Here $\Sigma_{ii}=1$ by spin-flip symmetry.  Positivity also controls the
cross block: for $u\in\R^G$ and $w\in\R^B$,
\[
 |u^T\Gamma_{GB}w|^2
 \le(u^T\Gamma_{GG}u)(w^T\Gamma_{BB}w).
\]
This follows by applying nonnegativity to the quadratic form at
$(u,sw)$ for all real $s$.  If
$a=\|\Gamma_{GG}\|$ and $b=\|\Gamma_{BB}\|$, it follows that
\[
 (u,w)^T\Gamma(u,w)
 \le(\sqrt a\,\|u\|+\sqrt b\,\|w\|)^2
 \le(a+b)(\|u\|^2+\|w\|^2).
\]
The same bound holds when either block is empty.  Therefore
$\|\Gamma\|\le16\|\Sigma\|+n^{13/25}$, proving
\eqref{eq:linear-covariance-comparison}.  Corollary~\ref{cor:covariance-scale}
and $13/25<2/3$ give the final assertion.
\end{proof}

\begin{proof}[Proof of Theorem~\ref{thm:linear-scale}]
The factor in the linear comparison is fixed, so it can be absorbed
while preserving every divergent threshold. Let $M_n\to\infty$ be
positive and deterministic. Apply Corollary~\ref{cor:covariance-scale}
with the divergent threshold
$M_n/32$, and intersect its event with that of
Proposition~\ref{prop:linear-scale}.  With probability tending to one,
\[
 \frac{32n^{2/3}}{M_n}\le\|\Sigma\|
 \le\frac{M_n n^{2/3}}{32},
 \qquad
 \|\Sigma\|\le\mathcal R_{\mathrm{lin}}(\mu)
 \le\frac{M_n n^{2/3}}2+n^{13/25}.
\]
Since $13/25<2/3$, the final term is at most
$M_nn^{2/3}/2$ for all sufficiently large $n$.  These inequalities
give the claimed common interval for both quantities.
\end{proof}

\subsection{Comparison with the high-temperature model}

For the upper bound we combine the high-probability norm estimate of
Proposition~\ref{prop:spectral} with the following fixed-temperature,
zero-field functional inequality. Its constant may depend on the chosen
temperature.

\begin{theorem}[High-temperature spectral gap
  {\cite[Theorem~1.1]{skgapcompanion}}]
\label{thm:high-temperature-gap}
Fix $0<\beta<1$.  Let $J$ be a symmetric matrix with zero diagonal and
independent entries $J_{ij}\sim N(0,\beta^2/n)$ for $i<j$, and let
\[
 \nu_J(x)=Z_J^{-1}\exp\bigl(\tfrac12x^TJx\bigr),
 \qquad x\in\{-1,1\}^n.
\]
There is a finite constant $C_\beta$, depending only on $\beta$, such that
\[
 \Prob_J\!\left\{
   \Var_{\nu_J}(f)\le C_\beta\mathcal D_{\nu_J}(f)
   \text{ for every }f:\{-1,1\}^n\to\R
 \right\}\longrightarrow1.
\]
Here $\mathcal D_{\nu_J}$ is the unscaled sum of conditional variances.
Consequently, the discrete-time heat-bath chain choosing one uniform
site has spectral gap at least $1/(nC_\beta)$ with probability tending
to one.
\end{theorem}

The complete proof is supplied by the separate companion
\cite[Theorem~1.1]{skgapcompanion}. The constant $C_\beta$ is required
only for each fixed $\beta<1$: we select $\beta$ after the desired
exponential tolerance and before taking the limit in $n$.

\begin{lemma}[Comparison of conditional-variance forms]
\label{lem:density-comparison}
Let $\rho,\nu$ be probability measures with full support on
$\{-1,1\}^n$.  Suppose
\[
 0<m\le\frac{\rho(x)}{\nu(x)}\le M<\infty
 \qquad\text{for every }x.
\]
If $\Var_\nu(f)\le C_\nu\mathcal D_\nu(f)$ for every $f$, then
\begin{equation}\label{eq:density-comparison}
 \Var_\rho(f)\le \frac{M}{m}C_\nu\mathcal D_\rho(f)
 \qquad\text{for every }f.
\end{equation}
In particular, for $\rho(x)\propto e^{H(x)}$ and
$\nu(x)\propto e^{\beta H(x)}$, with $\beta<1$, one may take
$M/m=e^{(1-\beta)\osc H}$.
\end{lemma}

\begin{proof}
The comparison works directly with least-squares characterizations;
one need not compare the conditional kernels pointwise. For any
full-support probability measure $\eta$,
\begin{equation}\label{eq:form-infima}
 \Var_\eta(f)=\inf_{c\in\R}\eta[(f-c)^2],
 \qquad
 \mathcal D_\eta(f)
 =\sum_i\inf_{h_i}\eta[(f-h_i(x_{-i}))^2],
\end{equation}
where every $h_i$ ranges independently over all functions of the other
coordinates.  For the second identity, conditional expectation gives
the orthogonal decomposition
\[
 \eta[(f-h_i)^2]
 =\eta[(f-\eta[f\mid x_{-i}])^2]
  +\eta[(\eta[f\mid x_{-i}]-h_i)^2].
\]
The first identity follows by the same decomposition with constants.
The density bounds and \eqref{eq:form-infima} imply
\[
 \Var_\rho(f)\le M\Var_\nu(f),
 \qquad
 \mathcal D_\rho(f)\ge m\mathcal D_\nu(f),
\]
which prove \eqref{eq:density-comparison}.  For the exponential
densities, $\rho(x)/\nu(x)$ is a constant times $e^{(1-\beta)H(x)}$;
the common normalization cancels in the ratio of its maximum and
minimum.
\end{proof}

\begin{proposition}[Subexponential upper bound]
\label{prop:subexponential-upper}
For every fixed $\epsilon>0$,
\[
 \Prob\left\{t_{\mathrm{mix}}\le e^{\epsilon n},\quad
                 k_{\mathrm{mix}}\le e^{\epsilon n}\right\}
 \longrightarrow1.
\]
\end{proposition}

\begin{proof}
We keep the order of parameters explicit: first the requested
exponential tolerance, then a fixed subcritical temperature, then the
dimension. Fix $\epsilon>0$ and choose $\beta\in(0,1)$ such that
$3(1-\beta)<\epsilon$.  Let
$\mu_\beta(x)\propto e^{\beta H_W(x)}$.
The zero-diagonal part of $\beta W$ satisfies exactly the hypotheses of
Theorem~\ref{thm:high-temperature-gap}; omitting its diagonal changes
the Hamiltonian by a constant.  With probability tending to one,
\[
 \Var_{\mu_\beta}(f)\le C_\beta\mathcal D_{\mu_\beta}(f)
 \qquad\text{for every }f.
\]
Intersect this event with $\|W\|\le3$, whose probability tends to one
by Proposition~\ref{prop:spectral}.  On their intersection,
\[
 \osc H_W\le3n,
 \qquad
 \mu_{\min}:=\min_x\mu(x)
 \ge2^{-n}e^{-\osc H_W}\ge e^{-(3+\log2)n}.
\]
Lemma~\ref{lem:density-comparison} now gives
\begin{equation}\label{eq:critical-poincare-upper}
 \Var_\mu(f)\le C_n\mathcal D_\mu(f),
 \qquad C_n=C_\beta e^{3(1-\beta)n}.
\end{equation}
We may increase $C_\beta$ to be at least one.

The comparison has produced a full Poincar\'e inequality. To obtain
worst-start total variation for both clocks, combine its spectral decay
with the lower bound on the smallest Gibbs mass. We give the conversion
explicitly so that the factor $n$ in attempted updates remains visible.  By Lemma~\ref{lem:heatbath-form},
\eqref{eq:critical-poincare-upper} gives a continuous-time gap at least
$1/C_n$ and a discrete-time gap at least $1/(nC_n)$.  Since $P$ is an
average of orthogonal projections, its spectrum is nonnegative.
Spectral decomposition on the mean-zero subspace therefore yields
\[
 \|e^{t\mathcal L}h\|_{L^2(\mu)}
 \le e^{-t/C_n}\|h\|_{L^2(\mu)},
 \qquad
 \|P^kh\|_{L^2(\mu)}
 \le e^{-k/(nC_n)}\|h\|_{L^2(\mu)}.
\]
For a starting configuration $x$, take
$h_x(y)=\mathbf1_{\{y=x\}}/\mu(x)-1$, whose squared norm is
$\mu(x)^{-1}-1$.  Reversibility identifies the transition density
relative to $\mu$, minus one, with $e^{t\mathcal L}h_x$ or $P^kh_x$.
The $L^1$ expression for total variation and Cauchy--Schwarz give
\begin{align}
 \max_x d_t^{\mathrm{ct}}(W,x)
 &\le\tfrac12\mu_{\min}^{-1/2}e^{-t/C_n},\label{eq:gap-tv-ct}\\
 \max_x d_k^{\mathrm{dt}}(W,x)
 &\le\tfrac12\mu_{\min}^{-1/2}e^{-k/(nC_n)}.
 \label{eq:gap-tv-dt}
\end{align}
With $B_n=\log2+\tfrac12(3+\log2)n$, these bounds imply
\[
 t_{\mathrm{mix}}\le C_nB_n,
 \qquad
 k_{\mathrm{mix}}\le\lceil nC_nB_n\rceil.
\]
The two right-hand sides grow at most as a fixed constant depending
on $\beta$ times $n^2e^{3(1-\beta)n}$, plus the rounding term.
Since $\beta$ is fixed and $3(1-\beta)<\epsilon$, both are at most
$e^{\epsilon n}$ for all sufficiently large $n$.
The two events used above have an intersection of probability tending
to one without any independence assumption.  Finally, the mixing
times depend only on the off-diagonal disorder, so averaging out the
auxiliary diagonal preserves the conclusion.
\end{proof}

\begin{proof}[Proof of Theorem~\ref{thm:mixing}]
The independent upper comparison is now proved. It remains to
extract worst-start lower bounds from the stronger realized-start
statement. Fix $\epsilon>0$ and apply Theorem~\ref{thm:gibbs-start} to
\[
 t_n=n^{2/3-\epsilon},\qquad
 k_n=\lfloor n^{5/3-\epsilon}\rfloor.
\]
For each clock, with probability tending to one there is a starting
configuration whose distance at the displayed time exceeds $1/4$.
Worst-start distance is nonincreasing in time by contraction of total
variation under a Markov kernel.  Hence
$t_{\mathrm{mix}}\ge t_n$ and $k_{\mathrm{mix}}>k_n$ on those events.
Since $k_{\mathrm{mix}}$ is an integer, the latter implies
$k_{\mathrm{mix}}\ge n^{5/3-\epsilon}$, including when $k_n=0$.
Intersect these two lower-bound events with the event in
Proposition~\ref{prop:subexponential-upper} to obtain all four asserted
bounds simultaneously.
\end{proof}

\begin{remark}
Theorem~\ref{thm:linear-scale} retains arbitrary divergent factors
because it concerns linear functions. The full Poincar\'e constant is
a supremum over all nonconstant functions, and its upper bound does not
follow by enlarging the class in that theorem. Proposition~\ref{prop:subexponential-upper}
is the independent all-function comparison proved here. The stronger
critical exponent results in \cite[Theorem~1.1]{criticalmixingcompanion}
use separate arguments and have fixed positive exponent slack in their
all-function upper bounds.
\end{remark}
